% ============================================================
\section{Variational-Discrete Realization and Crime Ledgers}
\label{sec:realization-calculus}
% ============================================================

\begin{quote}\small
\textbf{What this adds.}
The preceding sections determine exact visible laws, returned residuals,
workload-relative transfers, and coded closure data on the operator side. A
numerical host still has to realize those laws through weak forms, discrete
spaces, geometry maps, quadrature, stage splits, solver choices, and interface
transfers. This section makes that realization chain native.

\smallskip\noindent
\textbf{Boundary.}
This is a realization grammar, not a new foundation. It does not derive a
specific CFD, spectral-element, or CBS code from bedrock. It provides the
exact typed ledger in which any declared realization must record its departures
from exactness. Characteristic-based split schemes and high-order spectral
element hosts motivate this layer, but they enter only as external
realizations of the grammar developed here
\cite{NithiarasuCodinaZienkiewicz2006,Nek5000Theory2026,Nek5000CaseFiles2026,Nek5000Filters2026,NekNekTutorial2026}.
\end{quote}

\subsection{Variational laws and realization chains}

\begin{definition}[Variational law]
\label{def:variational-law}\lean{VariationalLaw}
A \emph{variational law} is a sextuple
\[
\mathcal V=(X,Y,a,\ell,\mathcal J,\mathcal I),
\]
where \(X\) and \(Y\) are finite-dimensional Hilbert spaces,
\[
a:X\times Y\to \bbR
\]
is a bilinear form,
\[
\ell:Y\to \bbR
\]
is a linear load functional,
\(\mathcal J\) is a declared family of linear observables on \(X\), and
\(\mathcal I\) is a declared family of linear invariants or constraints on
\(X\).
\end{definition}

\begin{definition}[Realization chain]
\label{def:realization-chain}\lean{RealizationChain}
Let \(\mathcal V=(X,Y,a,\ell,\mathcal J,\mathcal I)\) be a variational law. A
\emph{realization chain} at level \(h\) is a package
\[
\mathfrak R_h
=
(\mathcal V,X_h,Y_h,\Pi_h^X,\Pi_h^Y,\Phi_h,Q_h,\mathcal T_h,\mathcal S_h,\mathcal I_h)
\]
together with residual representatives
\[
R_h^{\mathrm{cont}},
\ R_h^{\mathrm{proj}},
\ R_h^{\mathrm{geom}},
\ R_h^{\mathrm{quad}},
\ R_h^{\mathrm{alias}},
\ R_h^{\mathrm{split}},
\ R_h^{\mathrm{lin}},
\ R_h^{\mathrm{alg}},
\ R_h^{\mathrm{bc}},
\ R_h^{\mathrm{iface}}
:X_h\times Y_h\to \bbR
\]
such that
\[
R_h^{\mathrm{cont}}(u_h,v_h)
:=
\ell(\Pi_h^Y v_h)-a(\Pi_h^X u_h,\Pi_h^Y v_h).
\]
The declared data have the following roles:
\begin{enumerate}[label=(\roman*), leftmargin=2em, itemsep=0.2em]
\item \(X_h,Y_h\) are the realized trial and test spaces.
\item \(\Pi_h^X,\Pi_h^Y\) map realized states and tests into the continuum
      law.
\item \(\Phi_h\) records geometry pullback/pushforward data.
\item \(Q_h\) records quadrature data.
\item \(\mathcal T_h\) records stage or time-splitting data.
\item \(\mathcal S_h\) records algebraic solver/preconditioner data.
\item \(\mathcal I_h\) records boundary and interface transfer data.
\end{enumerate}
The terminal residual
\[
R_h^{\mathrm{real}}:=R_h^{\mathrm{iface}}
\]
is the fully realized residual delivered by the declared host.
\end{definition}

\begin{definition}[Typed crime ledger]
\label{def:typed-crime-ledger}\lean{TypedCrimeLedger}\lean{typedCrimeLedger}
In the setting of Definition~\ref{def:realization-chain}, define the
\emph{typed realization crimes} by
\begin{align*}
\mathfrak C_h^{\mathrm{proj}} &:= R_h^{\mathrm{proj}}-R_h^{\mathrm{cont}},\\
\mathfrak C_h^{\mathrm{geom}} &:= R_h^{\mathrm{geom}}-R_h^{\mathrm{proj}},\\
\mathfrak C_h^{\mathrm{quad}} &:= R_h^{\mathrm{quad}}-R_h^{\mathrm{geom}},\\
\mathfrak C_h^{\mathrm{alias}} &:= R_h^{\mathrm{alias}}-R_h^{\mathrm{quad}},\\
\mathfrak C_h^{\mathrm{split}} &:= R_h^{\mathrm{split}}-R_h^{\mathrm{alias}},\\
\mathfrak C_h^{\mathrm{lin}} &:= R_h^{\mathrm{lin}}-R_h^{\mathrm{split}},\\
\mathfrak C_h^{\mathrm{alg}} &:= R_h^{\mathrm{alg}}-R_h^{\mathrm{lin}},\\
\mathfrak C_h^{\mathrm{bc}} &:= R_h^{\mathrm{bc}}-R_h^{\mathrm{alg}},\\
\mathfrak C_h^{\mathrm{iface}} &:= R_h^{\mathrm{iface}}-R_h^{\mathrm{bc}}.
\end{align*}
The tuple
\[
\mathfrak C_h
=
\bigl(
\mathfrak C_h^{\mathrm{proj}},
\mathfrak C_h^{\mathrm{geom}},
\mathfrak C_h^{\mathrm{quad}},
\mathfrak C_h^{\mathrm{alias}},
\mathfrak C_h^{\mathrm{split}},
\mathfrak C_h^{\mathrm{lin}},
\mathfrak C_h^{\mathrm{alg}},
\mathfrak C_h^{\mathrm{bc}},
\mathfrak C_h^{\mathrm{iface}}
\bigr)
\]
is the \emph{crime ledger} of the realization.
\end{definition}

\begin{theorem}[Exact residual ledger]
\label{thm:exact-residual-ledger}\lean{exact_residual_ledger}
For every realization chain,
\[
R_h^{\mathrm{real}}-R_h^{\mathrm{cont}}
=
\mathfrak C_h^{\mathrm{proj}}
+
\mathfrak C_h^{\mathrm{geom}}
+
\mathfrak C_h^{\mathrm{quad}}
+
\mathfrak C_h^{\mathrm{alias}}
+
\mathfrak C_h^{\mathrm{split}}
+
\mathfrak C_h^{\mathrm{lin}}
+
\mathfrak C_h^{\mathrm{alg}}
+
\mathfrak C_h^{\mathrm{bc}}
+
\mathfrak C_h^{\mathrm{iface}}
\]
pointwise on \(X_h\times Y_h\).
\end{theorem}

\begin{proof}
Substitute the definitions from Definition~\ref{def:typed-crime-ledger} and
sum. All intermediate residual representatives cancel telescopically:
\begin{align*}
\sum_{\alpha}\mathfrak C_h^{\alpha}
&=
\bigl(R_h^{\mathrm{proj}}-R_h^{\mathrm{cont}}\bigr)
+
\bigl(R_h^{\mathrm{geom}}-R_h^{\mathrm{proj}}\bigr)
+
\cdots
+
\bigl(R_h^{\mathrm{iface}}-R_h^{\mathrm{bc}}\bigr)\\
&=
R_h^{\mathrm{iface}}-R_h^{\mathrm{cont}}
=
R_h^{\mathrm{real}}-R_h^{\mathrm{cont}}.
\end{align*}
\end{proof}

\begin{corollary}[No hidden crime]
\label{cor:no-hidden-crime}\lean{no_hidden_crime}
If every typed crime in \(\mathfrak C_h\) vanishes, then
\[
R_h^{\mathrm{real}}=R_h^{\mathrm{cont}}.
\]
\end{corollary}

\begin{proof}
Apply Theorem~\ref{thm:exact-residual-ledger}.
\end{proof}

\begin{definition}[Realization exactness]
\label{def:realization-exactness}\lean{realizationExact}
A realization chain is \emph{realization-exact} if every entry of its crime
ledger vanishes identically.
\end{definition}

\subsection{Compatible mixed realization}

\begin{definition}[Fortin-compatible mixed realization]
\label{def:fortin-compatible-realization}\lean{FortinCompatibleRealization}
Let \(U\), \(U_h\), and \(M_h\) be finite-dimensional Hilbert spaces, and let
\[
b:U\times M_h\to \bbR,
\qquad
b_h:U_h\times M_h\to \bbR
\]
be continuous bilinear forms. A \emph{Fortin-compatible mixed realization}
consists of a linear lift
\[
F_h:U\to U_h
\]
and a constant \(C_F\ge 1\) such that
\[
b_h(F_hu,q_h)=b(u,q_h)
\qquad
(\forall u\in U,\ q_h\in M_h),
\]
and
\[
\|F_hu\|\le C_F\|u\|
\qquad
(\forall u\in U).
\]
\end{definition}

\begin{theorem}[Discrete inf-sup inherited from a Fortin lift]
\label{thm:discrete-infsup-from-fortin}\lean{discrete_infsup_from_fortin}
Assume the setting of Definition~\ref{def:fortin-compatible-realization}, and
suppose there is a constant \(\beta>0\) such that
\[
\sup_{u\in U\setminus\{0\}}
\frac{|b(u,q_h)|}{\|u\|}
\ge
\beta\|q_h\|
\qquad
(\forall q_h\in M_h).
\]
Then
\[
\sup_{u_h\in U_h\setminus\{0\}}
\frac{|b_h(u_h,q_h)|}{\|u_h\|}
\ge
\frac{\beta}{C_F}\|q_h\|
\qquad
(\forall q_h\in M_h).
\]
\end{theorem}

\begin{proof}
Fix \(q_h\in M_h\). By hypothesis, for every \(\varepsilon>0\) there exists
\(u\in U\setminus\{0\}\) such that
\[
|b(u,q_h)|
\ge
(\beta-\varepsilon)\|u\|\|q_h\|.
\]
Set \(u_h:=F_hu\). Then
\[
|b_h(u_h,q_h)|
=
|b(u,q_h)|
\ge
(\beta-\varepsilon)\|u\|\|q_h\|
\ge
\frac{\beta-\varepsilon}{C_F}\|u_h\|\|q_h\|.
\]
Taking the supremum over \(u_h\) and then letting \(\varepsilon\downarrow 0\)
gives the result.
\end{proof}

\begin{proposition}[Commuting constraint crime vanishes]
\label{prop:commuting-constraint-crime-vanishes}\lean{commuting_constraint_crime_vanishes}
Let \(B:U\to Z\), \(B_h:U_h\to Z_h\), and \(\Pi_h^Z:Z\to Z_h\) be linear maps.
If
\[
B_hF_h=\Pi_h^Z B,
\]
then the lifted constraint defect
\[
\mathfrak C_h^{\mathrm{compat}}(u):=B_hF_hu-\Pi_h^ZBu
\]
vanishes for every \(u\in U\).
\end{proposition}

\begin{proof}
This is exactly the stated commuting identity.
\end{proof}

\subsection{Geometry, quadrature, and aliasing}

\begin{definition}[Geometry-realized element form]
\label{def:geometry-realized-element-form}\lean{GeometryRealizedForm}
Let \(X_h\) and \(Y_h\) be realized spaces over an element family
\(\mathcal K_h\). For each \(K\in\mathcal K_h\), let
\[
\Phi_K:\widehat K\to K
\]
be a declared reference-to-physical map, and let
\[
\widehat g_{K,u_h,v_h}:\widehat K\to \bbR
\]
be the pulled-back integrand representing the local realized form. The
\emph{geometry-realized form} is
\[
a_h^{\mathrm{geom}}(u_h,v_h)
:=
\sum_{K\in\mathcal K_h}
\int_{\widehat K}\widehat g_{K,u_h,v_h}(\widehat x)\,d\widehat x.
\]
\end{definition}

\begin{definition}[Quadrature-realized form and crimes]
\label{def:quadrature-realized-form}\lean{QuadratureRealizedForm}
In the setting of Definition~\ref{def:geometry-realized-element-form}, let
\(Q_h=\{Q_K\}_{K\in\mathcal K_h}\) be elementwise quadrature rules and define
\[
a_h^Q(u_h,v_h)
:=
\sum_{K\in\mathcal K_h}Q_K[\widehat g_{K,u_h,v_h}].
\]
If \(a_h^{\mathrm{proj}}\) is a declared projected exact form on
\(X_h\times Y_h\), define the geometry and quadrature crimes by
\[
\mathfrak C_h^{\mathrm{geom}}(u_h,v_h)
:=
a_h^{\mathrm{geom}}(u_h,v_h)-a_h^{\mathrm{proj}}(u_h,v_h),
\]
\[
\mathfrak C_h^{\mathrm{quad}}(u_h,v_h)
:=
a_h^Q(u_h,v_h)-a_h^{\mathrm{geom}}(u_h,v_h).
\]
\end{definition}

\begin{theorem}[Geometry exactness kills geometry crime]
\label{thm:geometry-exactness-kills-geometry-crime}\lean{geometry_exactness_kills_geometry_crime}
If
\[
a_h^{\mathrm{geom}}(u_h,v_h)=a_h^{\mathrm{proj}}(u_h,v_h)
\qquad
(\forall u_h\in X_h,\ v_h\in Y_h),
\]
then \(\mathfrak C_h^{\mathrm{geom}}=0\).
\end{theorem}

\begin{proof}
This is the definition of \(\mathfrak C_h^{\mathrm{geom}}\).
\end{proof}

\begin{theorem}[Quadrature exactness kills quadrature crime]
\label{thm:quadrature-exactness-kills-quadrature-crime}\lean{quadrature_exactness_kills_quadrature_crime}
Assume that each element \(K\in\mathcal K_h\) carries a function class
\(\widehat{\mathfrak P}_K\) such that
\[
\widehat g_{K,u_h,v_h}\in \widehat{\mathfrak P}_K
\qquad
(\forall u_h\in X_h,\ v_h\in Y_h),
\]
and that
\[
Q_K[\varphi]=\int_{\widehat K}\varphi(\widehat x)\,d\widehat x
\qquad
(\forall \varphi\in \widehat{\mathfrak P}_K).
\]
Then \(\mathfrak C_h^{\mathrm{quad}}=0\).
\end{theorem}

\begin{proof}
For every \(u_h\in X_h\) and \(v_h\in Y_h\), the hypothesis gives
\[
Q_K[\widehat g_{K,u_h,v_h}]
=
\int_{\widehat K}\widehat g_{K,u_h,v_h}(\widehat x)\,d\widehat x
\qquad
(\forall K\in\mathcal K_h).
\]
Summing over elements yields \(a_h^Q(u_h,v_h)=a_h^{\mathrm{geom}}(u_h,v_h)\),
hence \(\mathfrak C_h^{\mathrm{quad}}(u_h,v_h)=0\).
\end{proof}

\begin{definition}[Paired aliasing package]
\label{def:paired-aliasing-package}\lean{PairedAliasingPackage}
Let \(S_h\subseteq X_h\) and \(T_h\subseteq Y_h\). The \emph{paired aliasing
package} associated with quadrature is
\[
\Omega_h^{\mathrm{alias}}=(S_h,T_h,B_h^{\mathrm{alias}}),
\]
where
\[
B_h^{\mathrm{alias}}(v_h,u_h):=\mathfrak C_h^{\mathrm{quad}}(u_h,v_h).
\]
\end{definition}

\begin{corollary}[Quadrature exactness kills workload-visible aliasing]
\label{cor:quadrature-exactness-kills-aliasing}\lean{quadrature_exactness_kills_aliasing}
If \(\mathfrak C_h^{\mathrm{quad}}=0\) on \(S_h\times T_h\), then the paired
aliasing package \(\Omega_h^{\mathrm{alias}}\) is zero on that workload, and
every one-sided observable transfer obtained from it by
Definition~\ref{def:paired-observable-package} also vanishes.
\end{corollary}

\begin{proof}
The first claim is exactly the definition of
\(B_h^{\mathrm{alias}}\). The induced one-sided observable transfers are
obtained by currying the same bilinear map, so they vanish as well.
\end{proof}

\begin{proposition}[Certified dealiasing on a selected workload]
\label{prop:certified-dealiasing-selected-workload}\lean{certified_dealiasing_selected_workload}
Let \(F_h:X_h\to X_h\) be linear. If
\[
\mathfrak C_h^{\mathrm{quad}}(F_hu_h,v_h)=0
\qquad
(\forall u_h\in S_h,\ v_h\in T_h),
\]
then the filtered workload \(F_hS_h\) is alias-free on the paired package
\(\Omega_h^{\mathrm{alias}}\).
\end{proposition}

\begin{proof}
This is the previous corollary applied to the filtered source family
\(F_hS_h\).
\end{proof}

\subsection{Stages, solvers, interfaces, and goals}

\begin{definition}[Stage bundle and split crime]
\label{def:stage-bundle-split-crime}\lean{StageBundle}\lean{stageBundleSplitCrime}
Let \(X_h\) be a realized state space. A \emph{stage bundle} at step
\(\Delta t\) is a composition
\[
\mathcal T_{\Delta t}
:=
\Psi^{(s)}_{\Delta t}\circ\cdots\circ\Psi^{(1)}_{\Delta t}
:X_h\to X_h.
\]
If \(\mathcal E_{\Delta t}:X_h\to X_h\) is a declared reference full-step
evolution, the \emph{split crime} is
\[
\mathfrak C_{\Delta t}^{\mathrm{split}}
:=
\mathcal T_{\Delta t}-\mathcal E_{\Delta t}.
\]
\end{definition}

\begin{proposition}[Stagewise invariants survive composition]
\label{prop:stagewise-invariants-survive-composition}\lean{stagewise_invariants_survive_composition}
Let \(I:X_h\to Z\) be any map. If
\[
I\circ \Psi^{(j)}_{\Delta t}=I
\qquad
(j=1,\dots,s),
\]
then
\[
I\circ\mathcal T_{\Delta t}=I.
\]
\end{proposition}

\begin{proof}
Apply the stage identities successively:
\[
I\circ\mathcal T_{\Delta t}
=
I\circ \Psi^{(s)}_{\Delta t}\circ\cdots\circ\Psi^{(1)}_{\Delta t}
=
I\circ \Psi^{(s-1)}_{\Delta t}\circ\cdots\circ\Psi^{(1)}_{\Delta t}
=
\cdots
=
I.
\]
\end{proof}

\begin{definition}[Solver package and algebraic crime]
\label{def:solver-package-algebraic-crime}\lean{SolverPackage}\lean{solverPackageAlgebraicCrime}
Let \(A_h:X_h\to X_h\) be a realized linear operator, let \(f_h\in X_h\), and
let \(\widetilde u_h\in X_h\) be a computed state. A \emph{solver package} is a
tuple
\[
\mathcal S_h=(A_h,M_h,\mathcal D_h,\mathcal P_h,\widetilde u_h),
\]
where \(M_h\) is a declared preconditioner, \(\mathcal D_h\) a declared
decomposition, and \(\mathcal P_h\) declared restriction/prolongation data. The
\emph{algebraic crime} is the residual
\[
r_h^{\mathrm{alg}}:=f_h-A_h\widetilde u_h.
\]
\end{definition}

\begin{theorem}[Exact solve iff algebraic crime vanishes]
\label{thm:exact-solve-iff-algebraic-crime-vanishes}\lean{exact_solve_iff_algebraic_crime_vanishes}
Assume \(A_h\) is invertible. Then
\[
r_h^{\mathrm{alg}}=0
\qquad\Longleftrightarrow\qquad
\widetilde u_h=A_h^{-1}f_h.
\]
\end{theorem}

\begin{proof}
If \(r_h^{\mathrm{alg}}=0\), then \(A_h\widetilde u_h=f_h\), so invertibility
forces \(\widetilde u_h=A_h^{-1}f_h\). The converse is immediate.
\end{proof}

\begin{proposition}[Spectral-equivalence solver certificate]
\label{prop:spectral-equivalence-solver-certificate}\lean{spectral_equivalence_solver_certificate}
Assume \(A_h\) and \(M_h\) are self-adjoint positive definite and that there
exist constants \(0<m\le M\) such that
\[
m\,\langle A_hx,x\rangle
\le
\langle M_hx,x\rangle
\le
M\,\langle A_hx,x\rangle
\qquad
(\forall x\in X_h).
\]
Then every eigenvalue \(\lambda\) of \(M_h^{-1}A_h\) lies in
\([
1/M,\,
1/m
]\),
so
\[
\kappa(M_h^{-1}A_h)\le M/m.
\]
\end{proposition}

\begin{proof}
Let \(A_hx=\lambda M_hx\) with \(x\ne 0\). Then
\[
\lambda
=
\frac{\langle A_hx,x\rangle}{\langle M_hx,x\rangle}.
\]
The spectral-equivalence bounds imply
\[
\frac{1}{M}
\le
\frac{\langle A_hx,x\rangle}{\langle M_hx,x\rangle}
\le
\frac{1}{m}.
\]
Thus \(\lambda\in[1/M,1/m]\), and the condition-number bound follows.
\end{proof}

\begin{definition}[Interface package and interface crime]
\label{def:interface-package-crime}\lean{InterfacePackage}\lean{interfaceCrime}
Let \(X_h\) and \(X_h^{\mathrm{ext}}\) be realized state spaces with trace maps
\[
\gamma_h:X_h\to \Lambda_h,
\qquad
\gamma_h^{\mathrm{ext}}:X_h^{\mathrm{ext}}\to \Lambda_h^{\mathrm{ext}}.
\]
An \emph{interface package} consists of a transfer map
\[
I_h:\Lambda_h^{\mathrm{ext}}\to \Lambda_h,
\]
a lift
\[
L_h:\Lambda_h\to X_h,
\]
and a declared host-transfer map
\[
T_h:X_h^{\mathrm{ext}}\to X_h.
\]
The \emph{interface crime} is
\[
\mathfrak C_h^{\mathrm{iface}}(u_h^{\mathrm{ext}})
:=
\gamma_hT_hu_h^{\mathrm{ext}}
-I_h\gamma_h^{\mathrm{ext}}u_h^{\mathrm{ext}}.
\]
\end{definition}

\begin{theorem}[Commuting interface diagrams kill interface crime]
\label{thm:commuting-interface-diagrams-kill-interface-crime}\lean{commuting_interface_diagrams_kill_interface_crime}
If
\[
\gamma_hT_h=I_h\gamma_h^{\mathrm{ext}},
\]
then \(\mathfrak C_h^{\mathrm{iface}}=0\). If moreover
\[
\gamma_hL_h=\operatorname{id}_{\Lambda_h},
\]
then the lifted interface correction \(L_h\mathfrak C_h^{\mathrm{iface}}\)
also vanishes.
\end{theorem}

\begin{proof}
The commuting identity is exactly the statement
\(\mathfrak C_h^{\mathrm{iface}}=0\). The second claim follows by applying
\(L_h\) to zero.
\end{proof}

\begin{definition}[Adjoint goal certificate]
\label{def:adjoint-goal-certificate}\lean{AdjointGoalCertificate}
Let \(A_h:X_h\to X_h\) be invertible and let \(\psi_h\in X_h\). The scalar goal
observable associated with \(\psi_h\) is
\[
J_{\psi_h}(v_h):=\langle \psi_h,v_h\rangle.
\]
The corresponding \emph{adjoint goal weight} is
\[
z_h:=(A_h^\ast)^{-1}\psi_h.
\]
\end{definition}

\begin{theorem}[Exact goal-error ledger identity]
\label{thm:exact-goal-error-ledger-identity}\lean{exact_goal_error_ledger_identity}
Let \(A_h:X_h\to X_h\) be invertible, let \(f_h\in X_h\), let
\(\widetilde u_h\in X_h\), and let
\[
u_h^\star:=A_h^{-1}f_h,
\qquad
e_h:=u_h^\star-\widetilde u_h,
\qquad
r_h:=f_h-A_h\widetilde u_h,
\]
and let \(z_h\) be the adjoint goal weight from
Definition~\ref{def:adjoint-goal-certificate}. If
\[
r_h=\sum_{\alpha} r_h^{\alpha},
\]
then
\[
J_{\psi_h}(e_h)
=
\langle z_h,r_h\rangle
=
\sum_{\alpha}\langle z_h,r_h^{\alpha}\rangle.
\]
\end{theorem}

\begin{proof}
Since \(A_hu_h^\star=f_h\),
\[
A_he_h
=
A_hu_h^\star-A_h\widetilde u_h
=
f_h-A_h\widetilde u_h
=
r_h.
\]
Therefore
\[
J_{\psi_h}(e_h)
=
\langle \psi_h,e_h\rangle
=
\langle A_h^\ast z_h,e_h\rangle
=
\langle z_h,A_he_h\rangle
=
\langle z_h,r_h\rangle.
\]
The final identity is linearity of the inner product.
\end{proof}

\begin{remark}[Realization certificates as tools]
\label{rem:realization-certificates-as-tools}\lean{exact_residual_ledger,discrete_infsup_from_fortin,geometry_exactness_kills_geometry_crime,quadrature_exactness_kills_quadrature_crime,stagewise_invariants_survive_composition,spectral_equivalence_solver_certificate,commuting_interface_diagrams_kill_interface_crime,exact_goal_error_ledger_identity}
The operator-side calculus already certifies visible/shadow splitting and
returned residuals. The present layer packages the complementary engineering
question: how a declared host realizes that law, and exactly where it does not.
The residual ledger, Fortin-compatible inheritance theorem, quadrature
exactness certificate, stagewise invariant theorem, spectral-equivalence
solver certificate, commuting-interface certificate, and dual-weighted
goal-error identity are the core theoremic outputs of this realization layer.
\end{remark}
